% Appendix C, grading sandbox and isolation (restructure plan, part 3). Sources: src/tracecontract/sandbox/ (the
% runner, the harness and its limits); tests/test_sandbox_isolation.py (the probes); docs/decision_log.md D28-D31, D35,
% D36, D40, D42, D44-D46, D51, D52, D54. Moved here: method.tex's process-separated grading, reproducibility.tex's
% isolation probes and exploit surface, limitations.tex's harness defects (condensed), the introduction's list of fixed
% forged-grade paths and the Harbor sentence. Every number is a macro (\dv: the harness limits, read from the code;
% \plv, \fv).
\section{Grading sandbox and isolation}
\label{app:grading}\label{sec:grading}
Hidden tests run in a trusted grader process (root) inside the container. Candidate code runs only in child processes
that a test starts through a small API (\texttt{run\_python}, \texttt{run\_cmd}), under an unprivileged user ID allocated
for that container alone, which cannot read the tests, write the results or signal the grader. A \texttt{run\_python}
child ignores user-site and startup hooks and imports the trusted modules that serialise and sign its result, those of
its error path included, before adding the candidate directories to its module path; when the snippet returns or
raises, it restores its own \texttt{sys.path} and \texttt{sys.meta\_path} and clears the importer cache before
formatting, serialising and signing. Code that the snippet imports still runs in that process and controls what that
call reports (below). A \texttt{run\_cmd} child starts with the candidate directories on \texttt{PYTHONPATH} and reports
nothing authenticated; no generated test, reference check or labeler check calls it, but the grader runs the
repository's tests through it.
\begin{itemize}
\item \emph{Authentication.} Each \texttt{run\_python} result is authenticated with a per-call key passed through an
  inherited pipe.
\item \emph{Trusted aggregation.} The host recomputes the outcome from a root-owned results file and a per-run nonce.
  Missing, stale, unparseable or incomplete results are verifier errors, never passes.
\item \emph{Container.} No network; all capabilities dropped except those needed to switch users, signal and hand over
  files; \texttt{no-new-privileges}, a read-only root filesystem, tmpfs scratch space and an init process.
\item \emph{Limits.} Each container gets \dv{harness.cpus} CPUs, \dv{harness.memory}\,GiB of memory and
  \dv{harness.pids} PIDs. Candidate processes also get CPU-time, \dv{harness.as}\,GiB address-space,
  \dv{harness.fsize}\,MiB file-size and \dv{harness.nproc}-process limits, with bounded malloc arenas.
\item \emph{Protected files.} Candidate edits to test files and test configuration (\texttt{tests/},
  \texttt{test\_*.py}, \texttt{conftest.py}, \texttt{pytest.ini}, \texttt{tox.ini}, \texttt{setup.cfg},
  \texttt{sitecustomize.py}, \texttt{*.pth}, \dots) are discarded before grading.
\item \emph{Patch policy.} Patches with path traversal, binary content or symbolic links are rejected.
\end{itemize}
Outcomes are \emph{pass}, \emph{task\_failure}, \emph{environment\_failure}, \emph{verifier\_error}, \emph{timeout} and
\emph{unresolved} (repeated runs disagree); only \emph{pass} yields reward.

\paragraph{Harness versions.} The pilot was graded under harness \plv{run.harness}; its re-grade and the fresh-agent run
use \fv{run.harness} (protocol v0.8), and so does the natural run (\nrv{run.harness}). That version differs from the
pilot's only in this: formatting
or encoding a value or an exception for a result message never turns a successful call into a failed call and never
leaves a call without a result. Values that JSON cannot encode are converted by trusted code that calls no candidate
method, or replaced by a placeholder naming their type; candidate code that runs during formatting can still end or
stall the call. Harness 0.3.4, whose fix was too narrow, was never used (D51, D52). Defects of earlier versions in
resource isolation (D29--D31), and paths through which a candidate module could forge signed results, which the
internal reviews found after the isolation probes had missed them (D36, D40, D44), were each fixed with a regression
probe before the runs reported here; runs under earlier versions are kept as superseded. In detail:
\begin{itemize}
\item The resource-isolation defects and the first start-up paths for forged results, the second left by the first fix,
  were found after evaluation runs and are fixed in harness 0.3.2.
\item A later verification found a further path for forged results in the error handling of 0.3.2, which harness 0.3.3
  closes; no candidate patch in the pool ships an \texttt{ast} or \texttt{unicodedata} module (D44). Every reported pilot
  result was re-run under 0.3.3, with the same outcome on every pair (D45).
\item In 0.3.3, methods that candidate objects define and that the wrapper calls while it formats a result
  (\texttt{\_\_repr\_\_}, \texttt{\_\_str\_\_}) run without the candidate directories on the module path, so such a method
  that imports a candidate module fails. No check in the pool is affected: all \plv{checks.results} check results equal those
  under 0.3.2 (D46).
\item A later verification found a second way in which 0.3.3 can report a successful call as failed: a result that JSON
  cannot encode, such as a dictionary with keys of mixed types (D51). Harness 0.3.5 keeps both from changing whether a
  call succeeded (D52), and a re-grade of all \fv{rg.pairs} pilot pairs under it gives the same outcome on every pair
  (\fv{rg.vs.pilotv033.outcome}; run \texttt{\REGRADERUN}, D54; Table~\ref{tab:regrade}).
\end{itemize}
The development audits ran under earlier harness versions.

\paragraph{Isolation probes.} The grading boundary is covered by Docker integration tests, run as part of the test
suite. Candidate code:
\begin{itemize}
\item cannot write the results directory, read the hidden tests, signal the grader or read the grader's environment;
\item has no network;
\item runs as an unprivileged user ID unique to the container;
\item is bounded by per-call timeouts, and its fork bombs are contained and daemonised processes reaped.
\end{itemize}
Forged or rewritten results are detected. Candidate edits to test files are removed before grading, and path-traversal
patches are rejected. Regression probes target the defects found during the pilot and its reviews:
\begin{itemize}
\item per-user limits are not shared with other containers;
\item concurrent UID allocations are distinct;
\item thread creation survives allocation churn under the address-space limit;
\item a candidate module that shadows the standard library cannot forge results;
\item a candidate \texttt{encodings} package, which the interpreter would import while it starts, cannot forge results;
\item startup hooks and packages written at runtime never load in later calls;
\item a candidate \texttt{ast} module, which Python's \texttt{traceback} imports while it formats an exception, cannot
  forge the report of a snippet that raises.
\end{itemize}
Further probes check harness 0.3.5: a candidate \texttt{\_\_repr\_\_} or \texttt{\_\_str\_\_} that fails while a result
is formatted becomes a placeholder, and a result that cannot be encoded never changes the outcome of a call.

\paragraph{Known exploit surface.} Within a sandbox call, candidate code controls the value it reports, and it can detect
that it is being graded. Tests therefore assert on concrete values computed by the trusted grader, but they cannot rule
out behaviour that differs under test. The repository-test baseline (B1) runs pytest with candidate code imported into
the reporting process, so a deliberately adversarial candidate can forge B1's report. A characterisation probe shows it:
a candidate whose exit hook rewrites the report passes, while an honestly wrong candidate fails (D42). B1 is meaningful
only for non-adversarial candidates, such as the held-out candidates here, and is unsuitable as a reward for
adversarial policies; the process-separated verifiers do not have this exposure, and they resist the tampering probes
above.

\paragraph{Harbor export.} The tasks' Harbor export uses a per-run sandbox user ID (D35); it can be re-validated from a
clone with \texttt{scripts/harbor\_validate.sh}, which records its outputs in \texttt{results/harbor\_validation/}.
